\chapter{Arquitectura del sistema}

La arquitectura se divide en siete capas. Primero se valida navegación y seguridad; después se añaden trazabilidad, integración y analítica por fases.

\begin{figure}[H]
\centering
\resizebox{\textwidth}{!}{%
\begin{tikzpicture}[
  x=1cm,y=1cm,
  block/.style={draw=UniPrimaryDeep!60,fill=white,rounded corners=2mm,align=left,font=\footnotesize,text width=2.45cm,minimum height=1.02cm,inner sep=4pt},
  core/.style={draw=UniPrimaryDeep,fill=UniPrimary!10,rounded corners=2mm,align=center,font=\footnotesize\bfseries,text width=2.60cm,minimum height=1.05cm,inner sep=4pt},
  safe/.style={draw=WarnAmber!85!black,fill=WarnAmber!10,rounded corners=2mm,align=left,font=\footnotesize,text width=2.45cm,minimum height=1.02cm,inner sep=4pt},
  arr/.style={-{Latex[length=2.3mm,width=1.6mm]},line width=0.75pt,draw=UniPrimaryDeep!75},
  dat/.style={<->,line width=0.7pt,draw=UniPrimary!80}
]

\fill[UniLight] (-0.15,-0.15) rectangle (12.85,5.25);

\node[block] (per) at (1.35,4.15) {\textbf{Percepci\'on}\\LiDAR, RGB-D, RFID/QR};
\node[block] (loc) at (4.35,4.15) {\textbf{Localizaci\'on}\\SLAM, IMU, encoders};
\node[core] (plan) at (7.35,4.15) {Planificador\\de misiones};
\node[block] (db) at (10.35,4.15) {\textbf{BD trazable}\\misiones, FOD, entregas};

\node[safe] (seg) at (1.35,2.20) {\textbf{Seguridad}\\esc\'aner, bumpers, E-stop};
\node[block] (ctrl) at (4.35,2.20) {\textbf{Control}\\base omnidireccional};
\node[block] (hmi) at (7.35,2.20) {\textbf{HMI y flota}\\avisos, bater\'ia, prioridad};
\node[block] (mes) at (10.35,2.20) {\textbf{MES/ERP}\\orden, kit y estaci\'on};

\draw[arr] (per) -- (loc);
\draw[arr] (loc) -- (plan);
\draw[dat] (plan) -- (db);
\draw[arr] (seg) -- (ctrl);
\draw[arr] (ctrl) -- (hmi);
\draw[dat] (hmi) -- (mes);
\draw[arr] (plan.south) -- (hmi.north);
\draw[arr] (seg.north) -- ++(0,0.45) -| (ctrl.north);
\draw[arr] (db.south) -- (mes.north);

\draw[WarnAmber!80!black,line width=0.95pt] (0.35,1.05) -- (11.40,1.05);
\node[anchor=west,font=\scriptsize\bfseries,text=WarnAmber!80!black] at (0.42,0.70)
{la capa de seguridad interrumpe cualquier misi\'on; la base de datos solo acepta misiones cerradas};
\end{tikzpicture}

}
\caption{Arquitectura propuesta para el Sistema AMR-FOD-Kitting.}
\label{fig:arquitectura-amr}
\end{figure}

\section{Capa de percepción}

La percepción combina sensores de navegación y sensores de proceso. El LiDAR proporciona detección geométrica y apoyo al mapa; la seguridad funcional queda separada en el escáner láser certificado, bumpers y parada de emergencia. La cámara RGB-D registra posibles FOD y documenta incidencias, pero sus positivos deben revisarse por un responsable de calidad. El sistema RFID/QR identifica bandejas, herramientas, órdenes y puntos de entrega.

\section{Planificación y supervisión}

El planificador asigna misiones según prioridad, estado del robot, proximidad y restricciones de zona. La HMI sirve para solicitar entregas, confirmar recepciones, cancelar misiones y consultar alarmas. La integración con MES/ERP empieza con importación y exportación de órdenes; en fases posteriores se habilita comunicación mediante API.

\section{Base de datos de trazabilidad}

Cada misión guarda identificador de kit, estación destino, hora de carga, hora de entrega, confirmación, ruta, incidencias, fotografía FOD si aplica y operador que valida. Con esos campos se audita el flujo y se miden cambios sin depender de impresiones subjetivas.
